\documentclass[11pt]{article}
\usepackage[margin=1in]{geometry}
\usepackage{amsmath,amssymb,amsthm}
\usepackage[T1]{fontenc}
\newcommand{\Z}{\mathbb{Z}}
\newcommand{\R}{\mathbb{R}}
\newcommand{\Q}{\mathbb{Q}}
\newtheorem{theorem}{Theorem}
\newtheorem{proposition}[theorem]{Proposition}
\newtheorem{corollary}[theorem]{Corollary}
\newtheorem{lemma}[theorem]{Lemma}
\theoremstyle{definition}
\newtheorem{remark}[theorem]{Remark}
\newtheorem{definition}[theorem]{Definition}

\title{Universality and first deviation for orthoscheme reflection groups\\ in every dimension}
\author{Vico Bonfioli}
\date{}

\begin{document}
\maketitle

\begin{abstract}
\noindent The right-triangle universality of the companion paper --- all shapes share a
growth prefix, then deviate at a shortest-vector law --- is the two-dimensional case of a
general dichotomy. We isolate the abstract mechanism: for any family of quotients of a
\emph{common} group, spherical growth is shape-independent through a horizon equal to half the
shortest added relator, and first deviates exactly there. We then exhibit the orthoscheme
(path-simplex) reflection groups as a family realizing this dichotomy in \emph{every} dimension:
generic shapes are rigid (all give one universal growth sequence), and every shape agrees with it
up to a finite, shape-dependent horizon. The universal sequences and their ratios
$r_n = 1+2\cos\frac{2\pi}{n+3}$ are computed for $n=2,3,4$; the arithmetic of the universal
sequence is dimension-dependent --- transcendental at $n=2$ (\textup{OEIS A396406}), rational at
$n=3$ ($9\cdot2^{\,d-2}$) --- so \emph{universality is general while transcendence is special to
the plane}.
\end{abstract}

\section{The abstract dichotomy}

Let $G=\langle S\rangle$ be a group with a finite symmetric generating set $S$, and let
$\ell_G$ denote word length. For a normal subgroup $N\trianglelefteq G$ the quotient
$Q=G/N$ inherits the generating set $\bar S$; write $s_d(H)=\#\{h\in H:\ell_H(h)=d\}$ for the
sphere sizes and $K(N)=\min\{\ell_G(g):g\in N\setminus\{1\}\}$ for the length of the shortest
nontrivial element of $N$.

\begin{theorem}[Universality--deviation dichotomy]
\label{thm:dichotomy}
Let $n:=\big\lceil K(N)/2\big\rceil$. Then
\[
   s_d(Q)=s_d(G)\quad\text{for all } d<n,
   \qquad\text{and}\qquad
   s_n(Q)<s_n(G).
\]
Equivalently, the first depth at which the growth of the quotient falls below that of $G$ is
exactly $\big\lceil K(N)/2\big\rceil$.
\end{theorem}

\begin{proof}
The projection $\pi:G\to Q$ is surjective and $1$-Lipschitz, so $\ell_Q(\pi g)\le\ell_G(g)$ and
$\pi$ maps the radius-$d$ ball $B_d(G)$ onto $B_d(Q)$. Thus $s_d(Q)\le s_d(G)$ for every $d$, with
equality iff $\pi$ is injective on the sphere $\Sigma_d(G)$ \emph{and} no two elements of
$\Sigma_d(G)$, $\Sigma_{d'}(G)$ with $d'\le d$ collapse together --- i.e.\ iff $\pi|_{B_d(G)}$ is
injective and length-preserving.

Suppose $d<n$, i.e.\ $2d<K(N)$. If $\pi g=\pi h$ with $g,h\in B_d(G)$, then $gh^{-1}\in N$ and
$\ell_G(gh^{-1})\le \ell_G(g)+\ell_G(h)\le 2d<K(N)$, forcing $gh^{-1}=1$, so $g=h$. Hence
$\pi|_{B_d(G)}$ is injective; it is also length-preserving, since $\ell_Q(\pi g)<\ell_G(g)$ would
give a shorter representative $g'$ with $\pi g'=\pi g$, whence $g(g')^{-1}\in N\setminus\{1\}$ has
length $\le 2d<K(N)$, again impossible. Therefore $s_d(Q)=s_d(G)$ for all $d<n$, and cumulatively
$|B_{n-1}(Q)|=|B_{n-1}(G)|$.

At depth $n$: choose $g_0\in N$ with $\ell_G(g_0)=K(N)$ and split a geodesic word for it as
$g_0=w_1w_2$ with $\ell_G(w_1)=\lceil K/2\rceil=n$ and $\ell_G(w_2)=\lfloor K/2\rfloor$. Both
subwords are geodesic (a shorter one would shorten $g_0$), and $\pi w_1=\pi(w_2^{-1})$ because
$w_1w_2\in N$; yet $w_1\ne w_2^{-1}$ in $G$ (else $g_0=1$). Thus two distinct elements of
$B_n(G)$ share a fiber, so $|B_n(Q)|<|B_n(G)|$; combined with $|B_{n-1}(Q)|=|B_{n-1}(G)|$ this
gives $s_n(Q)=|B_n(Q)|-|B_{n-1}(Q)|<s_n(G)$.
\end{proof}

\begin{remark}
The theorem is a statement about \emph{one} quotient. It becomes a \emph{universality} statement
when applied to a \emph{family} $\{Q_T=G/N_T\}$ sharing the common group $G$: every member has the
same growth as $G$ --- hence as every other member --- through its own horizon
$n_T=\lceil K(N_T)/2\rceil$, and members with larger $K(N_T)$ agree longer. The growth is
shape-independent exactly as far as the shortest added relator allows, and the first deviation
depth is a shortest-vector problem in $N_T$. This is the abstract content of the right-triangle
theorems of the companion paper, where $G$ is the generic (virtually lamplighter) group,
$N_T=2(t-1)\mu_T\,\Z[t^{\pm1}]$, and $n_T$ is the deviation-depth law.
\end{remark}

\section{The orthoscheme family in every dimension}

An \emph{$n$-orthoscheme} (path simplex) is a Euclidean simplex with vertices
$P_0,\dots,P_n$ whose consecutive edges $P_{k-1}P_k$ are mutually orthogonal; normalizing
$P_k=P_{k-1}+\ell_k e_k$ places edge $k$ along the axis $e_k$ with leg length $\ell_k>0$
($\ell_1=1$). Its $n+1$ facets carry reflections $R_0,\dots,R_n$ (the reflection $R_j$ in the
facet opposite $P_j$), generating a group $W_T\le\mathrm{Isom}(\R^n)$; for $n=2$ this is the
right-triangle side-reflection group of the companion paper. We measure the geodesic growth
$s_d(W_T)$: the number of distinct isometries of word length $d$ in $\{R_0,\dots,R_n\}$.

Two facets $R_i,R_j$ with $|i-j|\ge2$ are orthogonal, hence commute; consecutive facets meet at a
dihedral angle that, for generic legs, is an irrational multiple of $\pi$. As with right
triangles, this makes the \emph{generic} orthoscheme group rigid: its isomorphism type, and hence
its growth, is independent of the (generic) shape. Computation confirms this and pins the
universal sequences.

\begin{proposition}[Generic universality, computational]
\label{prop:generic}
For $n\in\{2,3,4\}$, all tested generic $n$-orthoschemes share one geodesic-growth sequence
$U^{(n)}$, and the successive ratios converge to $r_n=1+2\cos\frac{2\pi}{n+3}$:
\[
\begin{array}{lll}
n=2: & U^{(2)}=1,3,5,8,13,21,34,55,89,144,225,\dots & (\textup{OEIS A396406})\\[2pt]
n=3: & U^{(3)}=1,4,9,18,36,72,144,288,\dots = \{1,4\}\ \text{then}\ 9\cdot2^{\,d-2} & (r_3=2,\ \text{exact to depth }16)\\[2pt]
n=4: & U^{(4)}=1,5,14,33,75,169,380,854,1919,4312,9689,\dots & (r_4=1+2\cos\tfrac{2\pi}7=2.24698\dots)
\end{array}
\]
The agreement is over four generic shapes per dimension, to depth $14$ ($n=2$), $16$ ($n=3$),
$12$ ($n=4$); the ratio $r_3=2$ and $r_4=2.24698$ are reproduced to full precision. (Script
\texttt{code/zeta\_probe/orthoscheme\_universality.py}.)
\end{proposition}

\begin{proposition}[Finite-horizon deviation, exact]
\label{prop:deviation}
Special shapes agree with $U^{(n)}$ up to a shape-dependent depth and then deviate. In dimension
$3$, in \emph{exact rational arithmetic}:
\[
\begin{array}{llll}
(1,1,1): & 1,4,9,\mathbf{17},\dots & \text{deviates at depth }3; &
(1,1,2),(2,1,1): 1,4,9,18,\mathbf{35},\dots\ \text{depth }4;\\
(1,\sqrt3,1): & \text{depth }6; & (1,2,3): \text{agrees past depth }14 & (\text{deep deviation}).
\end{array}
\]
The deviations are certified in exact arithmetic (rational legs give rational reflections), so they
are genuine group-theoretic collapses, not numerical artifacts. As in the plane, shapes with a
rational dihedral angle deviate early, while shapes with only ``generic'' angles (e.g.\ $(1,2,3)$)
agree far past the reach of direct enumeration.
\end{proposition}

\section{What is proven, and the significance}

Theorem~\ref{thm:dichotomy} is unconditional and dimension-free: it reduces the entire
universality-then-deviation phenomenon to two group-theoretic inputs --- a common ambient group
$G$, and the shortest-relator length $K(N_T)$ of each shape's kernel. For right triangles both
inputs are supplied in closed form by the companion paper (the lamplighter identification and the
ideal $2(t-1)\mu_T\Z[t^{\pm1}]$), making the $n=2$ instance a theorem outright, deviation law and
all. Propositions~\ref{prop:generic}--\ref{prop:deviation} establish the two inputs
\emph{computationally} in dimensions $3$ and $4$: the generic group is rigid (a universal sequence
exists), and each shape's horizon is finite. What remains open beyond $n=2$ is the structural
identification of the generic orthoscheme group $G_n$ and of the kernels $N_T$ --- the higher
analogue of the lamplighter theorem --- which would upgrade Propositions~\ref{prop:generic} and
\ref{prop:deviation} to theorems and yield a closed-form horizon $n_T$ in every dimension.

The phenomenon is thus \emph{general}: a whole family of reflection groups, in every dimension,
grows identically until an arithmetic obstruction of the individual shape switches on. Against this
backdrop the plane is special in one precise respect --- \emph{arithmetic}. The universal sequence
is transcendental only at $n=2$ (A396406, growth rate the transcendental $\beta_2=1.4916\dots$,
strictly below $r_2=\varphi$); at $n=3$ it is the rational $9\cdot2^{\,d-2}$, and the ratios
$r_3=2,\ r_4=2.24698\dots$ govern cleanly. Universality is the rule in every dimension;
transcendence is the exception, confined to the plane.

\end{document}
